%% Supplementary material for the DEFT paper
%% "Causal Models as Database Objects: The Case for a Causal DBMS".
%% Compile from this folder with XeLaTeX or LuaLaTeX:
%%   cd docs && latexmk -xelatex supplementary.tex
%% Listings are read from the repository root (inputpath=..).
\documentclass[a4paper,11pt]{article}
\usepackage[margin=2.3cm]{geometry}
\usepackage{fontspec}
\IfFontExistsTF{Menlo}{\setmonofont{Menlo}[Scale=0.85]}{\setmonofont{DejaVu Sans Mono}[Scale=0.85]}
\usepackage{listings,xcolor}
\usepackage{booktabs,array}
\usepackage[hidelinks]{hyperref}

\lstdefinelanguage{causal}{
  morekeywords={DEFINE,CAUSAL,VARIABLE,DAG,MODEL,FOR,UNIT,OBSERVED,AS,TYPE,ROLE,MEAN,NOISE,WEIGHT,
    CAUSES,EFFECTS,OF,ANCESTORS,DESCENDANTS,MEDIATORS,COLLIDERS,CONFOUNDERS,BETWEEN,AND,BACKDOOR,PATHS,
    IDENTIFY,EFFECT,ON,RETURN,ADJUSTMENT,SETS,ESTIMATE,USING,INTERVENE,COMPARE,TO,FROM,DISTRIBUTION,
    EXPECTATION,COUNTERFACTUAL,USE,RESET,SHOW,MODELS,DESCRIBE},
  sensitive=true, morecomment=[l]{//}, morecomment=[s]{/*}{*/}, morestring=[b]"}
\lstset{basicstyle=\ttfamily, keywordstyle=\color{blue!60!black}\bfseries,
  commentstyle=\color{gray}, stringstyle=\color{red!50!black},
  breaklines=true, frame=single, columns=fullflexible,
  literate={→}{{$\rightarrow$}}1 {≠}{{$\neq$}}1 {⊥}{{$\perp$}}1,
  inputpath=..}

\title{Supplementary Material for ``Causal Models as Database Objects:\\
The Case for a Causal DBMS''}
\author{Amedeo Pachera, Kirollos Nayrouz, Angela Bonifati, Andrea Mauri\\[2pt]
\normalsize CNRS LIRIS, Universit\'e Claude Bernard Lyon 1, Lyon, France}
\date{}

\begin{document}
\maketitle

\noindent This document complements the paper ``Causal Models as Database
Objects: The Case for a Causal DBMS''. It presents the declarative causal
query language that we propose for a Y-DBMS, whose statements compile to the
procedures of Section~IV of the paper. Section~\ref{sec:grammar} gives the
grammar of the language, Section~\ref{sec:example} shows it on the running
example of the paper, and Section~\ref{sec:dag} lists the variables and the
edges of that example. The code of the prototype, of the language, and of the
experiments is available at \url{https://github.com/pake97/Y-DBMS}.

\section{Grammar}\label{sec:grammar}
\lstinputlisting{dsl/src/main/antlr4/com/example/causal/dsl/parser/CausalDSL.g4}

\section{Example}\label{sec:example}
\lstinputlisting[language=causal]{dev/example.causal}

\section{The running example}\label{sec:dag}

Table~\ref{tab:dag} lists the $22$ variables of the policy-making DAG of
Fig.~2 of the paper, the organization that holds each of them, and their
parents, which give the $37$ edges of the DAG. Latent variables are held by
no organization. The treatment of the policy question is
\textit{energy\_subsidy} and the outcome is \textit{loan\_defaults}.

\begin{table}[h]\centering\small
\caption{Variables of the policy-making DAG, with their holder and their parents.}
\label{tab:dag}
\begin{tabular}{@{}l >{\raggedright\arraybackslash}p{0.2\linewidth} >{\raggedright\arraybackslash}p{0.42\linewidth}@{}}
\toprule
Variable & Holder & Parents \\
\midrule
\textit{regional\_unemployment} & statistics office & -- \\
\textit{household\_size} & statistics office & -- \\
\textit{urbanicity} & statistics office & -- \\
\textit{job\_loss} & employment and social services & \textit{regional\_unemployment}, \textit{local\_economic\_shock}, \textit{health\_shock} \\
\textit{employment\_status} & employment and social services & \textit{job\_loss} \\
\textit{benefit\_uptake} & employment and social services & \textit{household\_size}, \textit{financial\_literacy} \\
\textit{household\_income} & tax office & \textit{employment\_status} \\
\textit{eligibility\_rule} & policy authority & -- \\
\textit{energy\_subsidy} ($T$) & policy authority & \textit{eligibility\_rule}, \textit{energy\_price}, \textit{household\_income}, \textit{benefit\_uptake} \\
\textit{energy\_price} & energy providers & -- \\
\textit{energy\_consumption} & energy providers & \textit{household\_size}, \textit{urbanicity}, \textit{energy\_price} \\
\textit{energy\_arrears} & energy providers & \textit{energy\_consumption}, \textit{energy\_price}, \textit{energy\_subsidy} \\
\textit{credit\_access} & banks & \textit{urbanicity}, \textit{energy\_arrears} \\
\textit{consumer\_credit} & banks & \textit{credit\_access}, \textit{household\_income}, \textit{energy\_subsidy}, \textit{informal\_income} \\
\textit{savings\_buffer} & banks & \textit{household\_income}, \textit{financial\_literacy} \\
\textit{loan\_defaults} ($Y$) & banks & \textit{consumer\_credit}, \textit{savings\_buffer}, \textit{energy\_arrears}, \textit{household\_income}, \textit{regional\_unemployment}, \textit{insurance\_coverage}, \textit{health\_shock}, \textit{local\_economic\_shock}, \textit{informal\_income} \\
\textit{insurance\_coverage} & insurers & \textit{financial\_literacy} \\
\textit{insurance\_claims} & insurers & \textit{health\_shock}, \textit{insurance\_coverage} \\
\textit{health\_shock} & health service & -- \\
\textit{local\_economic\_shock} & latent & -- \\
\textit{financial\_literacy} & latent & -- \\
\textit{informal\_income} & latent & -- \\
\bottomrule
\end{tabular}
\end{table}

\end{document}
